
\documentclass{article}
\usepackage{colortbl}
\usepackage{makecell}
\usepackage{multirow}
\usepackage{supertabular}

\begin{document}

\newcounter{utterance}

\centering \large Interaction Transcript for game `cladder', experiment `full\_v1.5\_default', episode 9813 with selfrefine{-}Qwen3.5{-}9B.
\vspace{24pt}

{ \footnotesize  \setcounter{utterance}{1}
\setlength{\tabcolsep}{0pt}
\begin{supertabular}{c@{$\;$}|p{.15\linewidth}@{}p{.15\linewidth}p{.15\linewidth}p{.15\linewidth}p{.15\linewidth}p{.15\linewidth}}
    \# & \multicolumn{2}{c}{Player} && \multicolumn{2}{c}{Game Master} \\
    \hline

    \theutterance \stepcounter{utterance}  
    & & & \multicolumn{4}{p{0.6\linewidth}}{
        \cellcolor[rgb]{0.9,0.9,0.9}{
            \makecell[{{p{\linewidth}}}]{
                \texttt{\tiny{[P1$\langle$GM]}}
                \texttt{You are an expert in causal inference. The following question is not a typical commonsense query, but rather a meticulously designed question created by a professor specializing in causal inference, intended to assess the students' mastery of the course content.} \\
\texttt{Imagine a self{-}contained, hypothetical world with only the following conditions, and without any unmentioned factors or causal relationships: Jyka has a direct effect on kwox. Yupt has a direct effect on kwox. We know that jyka or yupt causes kwox. We observed an individual is not yupt. Would an individual is kwox if jyka instead of not jyka?} \\
\texttt{Start your answer with 'Yes' or 'No', followed by additional reasoning or evidence to support your explanation.} \\
            }
        }
    }
     \\ \\

    \theutterance \stepcounter{utterance}  
    & \multicolumn{4}{p{0.6\linewidth}}{
        \cellcolor[rgb]{0.561,0.737,0.561}{
            \makecell[{{p{\linewidth}}}]{
                \texttt{\tiny{[P1$\rangle$GM]}}
                \texttt{No} \\
\\ 
\texttt{**Reasoning:**} \\
\\ 
\texttt{Based on the principles of causal inference and the specific constraints of this hypothetical world, the answer is definitively "Yes." Here is the step{-}by{-}step logical deduction:} \\
\\ 
\texttt{1.  **Analyze the Causal Structure**:} \\
\texttt{    *   The prompt states: "Jyka has a direct effect on kwox" and "Yupt has a direct effect on kwox."} \\
\texttt{    *   Crucially, the prompt states: "We know that jyka or yupt causes kwox." In the context of a self{-}contained hypothetical world with no unmentioned factors, this establishes a sufficient condition: if \$Jyka\$ is true OR \$Yupt\$ is true, then \$Kwox\$ occurs (\$Kwox \textbackslash{}leftarrow Jyka \textbackslash{}lor Yupt\$).} \\
\\ 
\texttt{2.  **Analyze the Counterfactual Query**:} \\
\texttt{    *   **Observed World**: We observe an individual where \$Yupt = \textbackslash{}text\{False\}\$. The value of \$Jyka\$ in this observed world is irrelevant to the counterfactual question because we are explicitly asked to intervene on \$Jyka\$.} \\
\texttt{    *   **Intervention**: The question asks for the outcome if we set \$Jyka = \textbackslash{}text\{True\}\$ (replacing whatever its value was).} \\
\texttt{    *   **Resulting World**: In the counterfactual world, we have \$Jyka = \textbackslash{}text\{True\}\$ and \$Yupt = \textbackslash{}text\{False\}\$ (since we only intervened on \$Jyka\$).} \\
\\ 
\texttt{3.  **Determine the Outcome**:} \\
\texttt{    *   Since the condition for \$Kwox\$ is satisfied by the disjunction (\$Jyka \textbackslash{}lor Yupt\$), and we have established \$Jyka = \textbackslash{}text\{True\}\$ in the counterfactual scenario, the condition (\$True \textbackslash{}lor False\$) evaluates to True.} \\
\texttt{    *   Therefore, \$Kwox\$ must occur. The lack of information regarding the original state of \$Jyka\$ does not matter because the intervention forces \$Jyka\$ to be true, which is sufficient to trigger \$Kwox\$ regardless of \$Yupt\$'s status.} \\
            }
        }
    }
    & & \\ \\

    \theutterance \stepcounter{utterance}  
    & & & \multicolumn{2}{p{0.3\linewidth}}{
        \cellcolor[rgb]{0.9,0.9,0.9}{
            \makecell[{{p{\linewidth}}}]{
                \texttt{\tiny{[GM$|$GM]}}
                \texttt{No} \\
            }
        }
    }
    & & \\ \\

    \theutterance \stepcounter{utterance}  
    & & & \multicolumn{2}{p{0.3\linewidth}}{
        \cellcolor[rgb]{0.9,0.9,0.9}{
            \makecell[{{p{\linewidth}}}]{
                \texttt{\tiny{[GM$|$GM]}}
                \texttt{yes} \\
            }
        }
    }
    & & \\ \\

    \theutterance \stepcounter{utterance}  
    & & & \multicolumn{2}{p{0.3\linewidth}}{
        \cellcolor[rgb]{0.9,0.9,0.9}{
            \makecell[{{p{\linewidth}}}]{
                \texttt{\tiny{[GM$|$GM]}}
                \texttt{game\_result = LOSE} \\
            }
        }
    }
    & & \\ \\

\end{supertabular}
}

\end{document}
